\documentclass[11pt,a4paper]{amsart}

\usepackage[margin=2.3cm]{geometry}
\usepackage{amsmath,amssymb,amsthm}
\usepackage{booktabs}
\usepackage{xcolor}
\usepackage{hyperref}
\usepackage{microtype}
\sloppy

\hypersetup{
  colorlinks=true,
  linkcolor=blue!60!black,
  citecolor=blue!60!black,
  urlcolor=blue!60!black,
  pdftitle={A Machine-Certified Degree-2,799 Counterexample to Smale's Mean Value Conjecture via Envelope-Theorem Barycentric Minimax Optimization and pandroots},
  pdfauthor={Ivan Besevic}
}

\newtheorem{theorem}{Theorem}
\newtheorem{corollary}[theorem]{Corollary}
\theoremstyle{remark}
\newtheorem{remark}[theorem]{Remark}

\newcommand{\C}{\mathbb{C}}
\newcommand{\Q}{\mathbb{Q}}
\newcommand{\R}{\mathbb{R}}
\newcommand{\Z}{\mathbb{Z}}
\DeclareMathOperator{\RePart}{Re}

\title[A Degree-2,799 Counterexample to Smale's Mean Value Conjecture]{%
A Machine-Certified Degree-2,799 Counterexample to Smale's Mean Value Conjecture
via Envelope-Theorem Barycentric Minimax Optimization and \texttt{pandroots}}

\author{Ivan Besevic}
\address{ESIEE-IT, France}
\email{ivan.besevic@edu.esiee-it.fr}
\date{November 25, 2026}

\subjclass[2020]{Primary 30C10, 30C15; Secondary 65E05, 65G20, 68V15}
\keywords{Smale's mean value conjecture, critical points of polynomials,
Fej\'er--Lagrange interpolation, minimax optimization, Rouch\'e's theorem,
certified root finding, pandroots}

\begin{document}

\begin{abstract}
We construct and unconditionally machine-certify an explicit complex polynomial
$p(u) \in \Q(i)[u]$ of exact degree $d = 2{,}799$ satisfying $p(0) = 0$,
$p'(0) = 1$, and
\[
  \min_{p'(\zeta) = 0} \left|\frac{p(\zeta)}{\zeta}\right|
  \in \bigl[\,1.0002520813917934477351916,\;\;
              1.0002520813917934477351917\,\bigr] > 1,
\]
setting a new lowest-degree record for explicit counterexamples to Smale's 1981
Mean Value Conjecture ($K \le 1$). Compared with the October 2026 disproof of
Adamczewski and Colbrook (\texttt{arXiv:2610.36158}, degree $29{,}984$, minimum
$\approx 1.000000495958$) and our November 24, 2026 record (degree $7{,}299$,
minimum $\approx 1.000027196554$), the new construction reduces the certified
degree by $10.71\times$ and $2.61\times$ respectively, while increasing the
certified margin above $1$ by $508.3\times$ and $9.27\times$. By standard
leading-term perturbation, counterexamples exist for every degree
$n \ge 2{,}799$.

To break through the $\sim 7{,}300$ degree floor of pure analytic conformal-map
interpolation, we introduce two algorithmic techniques: (i) Envelope-Theorem
barycentric minimax optimization directly in the $O(\log M)$-conditioned
Fej\'er--Lagrange nodal basis with a free anchor $a \in \C$, staged from low-mode
Fourier saddle crossing to full nodal L-BFGS equalization; and (ii) iterative
bootstrap degree compression ($3{,}999 \to 3{,}499 \to 3{,}199 \to 2{,}799$)
coupled with Maehly's deflated Newton root tracking. The resulting compact
$1.3\text{ MB}$ exact Gaussian-integer polynomial is certified with zero
floating-point assumptions in $10.43\text{ s}$ on 10 CPU cores by combining
\texttt{pandroots v0.31.0} (\texttt{certify\_roots\_exact}, certifying all
$2{,}798$ critical points in $8.56\text{ s}$ to radius $\le 2^{-795}$) with an
exact arbitrary-precision integer Smale enclosure ($1.85\text{ s}$), and is
independently cross-checked at $700$-digit precision in Python \texttt{mpmath}.
\end{abstract}

\maketitle

\section{Introduction and Main Certified Theorem}

In 1981, Smale~\cite{Smale81} proved that every complex polynomial $p$ of degree
$d \ge 2$ normalized by $p(0) = 0$ and $p'(0) = 1$ possesses at least one
critical point $\zeta \in \C$ ($p'(\zeta) = 0$) satisfying
$|p(\zeta)/\zeta| \le 4$, and conjectured that the bound $4$ could be improved
to $1$ (or $(d-1)/d$). Following forty-five years of partial results establishing
the conjecture for small degrees and special classes, the global $K \le 1$
conjecture was refuted in October 2026 by Adamczewski and Colbrook~\cite{AC26}
via a certified counterexample of degree $29{,}984$, and Goel~\cite{Goel2026}
subsequently established that the global supremum over all degrees is sharp at
$K^* = 4$. On November 24, 2026, using a Fourier-smoothed analytic Pac-Man
domain and Fej\'er--Lagrange interpolation, we lowered the lowest explicitly
certified counterexample degree to $d = 7{,}299$~\cite{Besevic7299}.

A natural quantitative question remains open: \emph{what is the smallest
polynomial degree $d$ for which a counterexample to Smale's Mean Value Conjecture
exists?} Here we shatter the degree-$7{,}299$ threshold by a factor of $2.61\times$
by moving beyond pure conformal-map interpolation to direct nonlinear minimax
optimization over the full polynomial space.

\begin{theorem}[Machine-Certified Degree-2,799 Counterexample]\label{thm:smale2799}
Let $P(z) = \sum_{k=0}^{2799} c_k z^k \in \Z[i][z]$ be the explicit
Gaussian-integer polynomial archived in \texttt{data/champion\_deg2799.gpoly}
(SHA-256: \texttt{\scriptsize 5d1ad7a2509a488934031144d6c2bee5ab36d027bcee392087657a1b117dbfdc},
leading coefficient $c_{2799} \neq 0$, maximum coefficient bit-length $1{,}011$,
leading coefficient bit-length $558$), and let
\[
  a = \frac{22901 - 11056\,i}{2^{15}} \in \Q(i)
\]
be the exact $15$-bit dyadic anchor. Define the normalized polynomial
\[
  p(u) = \frac{P(u + a) - P(a)}{P'(a)} \in \Q(i)[u],
  \qquad \deg(p) = 2{,}799, \quad p(0) = 0, \quad p'(0) = 1.
\]
Then all $2{,}798$ critical points of $P$ in $\C$ are simple, lie strictly inside
$|c| < 96/100$, and satisfy the unconditionally certified global enclosure
\begin{equation}\label{eq:certified-min}
\begin{aligned}
  \min_{p'(\zeta) = 0} \left|\frac{p(\zeta)}{\zeta}\right|
  &= \min_{P'(c) = 0} \left|\frac{P(c) - P(a)}{(c - a)\, P'(a)}\right| \\
  &\in \bigl[\,1.0002520813917934477351916,\;\;
               1.0002520813917934477351917\,\bigr] > 1.
\end{aligned}
\end{equation}
Furthermore, all $2{,}798$ critical value ratios satisfy
$|p(\zeta)/\zeta| \le 1.10226259$.
\end{theorem}

\begin{corollary}
By adding $\varepsilon u^n$ with sufficiently small $\varepsilon > 0$ (via
Rouch\'e's theorem, as in~\cite{AC26}), counterexamples to Smale's Mean Value
Conjecture ($K \le 1$) exist for every polynomial degree $n \ge 2{,}799$.
\end{corollary}

Table~\ref{tab:comparison} summarizes the progression of rigorously certified
explicit counterexamples.

\begin{table}[h]
\centering
\footnotesize
\setlength{\tabcolsep}{4pt}
\caption{Progression of machine-certified explicit counterexamples to Smale's
Mean Value Conjecture ($K \le 1$).}
\label{tab:comparison}
\begin{tabular}{@{}lccc@{}}
\toprule
Metric & Adamczewski--Colbrook~\cite{AC26} & Besevic~\cite{Besevic7299} & \textbf{This Work} \\
       & (Oct 2026) & (Nov 24, 2026) & \textbf{(Nov 25, 2026)} \\
\midrule
Exact degree $d$ & $29{,}984$ & $7{,}299$ & $\mathbf{2{,}799}$ ($2.61\times$ / $10.71\times$) \\
Certified min & $(1.0000004959\ldots 01,\ldots 04)$ & $[1.0000271965\ldots 59, \ldots 60]$ & $\mathbf{[1.0002520813\ldots 16, \ldots 17]}$ \\
Margin above $1$ & $4.96 \times 10^{-7}$ & $2.72 \times 10^{-5}$ & $\mathbf{2.52 \times 10^{-4}}$ ($9.27\times$ / $508\times$) \\
Exact anchor $a$ & --- & $53$-bit dyadic & $\mathbf{(22901 - 11056\,i)/2^{15}}$ \\
Artifact size & --- & $7.6\text{ MB}$ ($2{,}311$ bits) & $\mathbf{1.3\text{ MB}}$ ($1{,}011$ bits) \\
Root certifier & --- & Custom BigInt Rouch\'e & \textbf{\texttt{pandroots v0.31.0}}~\cite{pandroots} \\
Total cert.\ time & --- & $413.72\text{ s}$ (10 cores) & $\mathbf{10.43\text{ s}}$ ($39.7\times$ faster) \\
\bottomrule
\end{tabular}
\end{table}

\section{Breaking the Analytic Interpolation Floor}

\subsection{Why Conformal-Map Interpolation Plateaus Near Degree 7,300}
In both~\cite{AC26} and~\cite{Besevic7299}, the polynomial $P$ is obtained as a
linear polynomial approximation (Fej\'er--Lagrange or Faber interpolant) to a
fixed univalent conformal map $f : \Omega \to \mathbb{D}$ on a slotted domain
$\Omega$. While the analytic Pac-Man boundary of~\cite{Besevic7299}
($h = 0.14$, $x_c = 0.32$, $\sigma = 5.0$, $k_{\max} = 28$) achieves an interior
conformal margin of $1.002027$ at $a = 0.72 - 0.38\,i$, its Fej\'er--Lagrange
interpolant of degree $d = M - 1$ suffers severe non-uniformity: for
$M < 7{,}000$, critical points on the outer arc have substantial surplus margin
($1.005$--$1.04$), whereas a few hundred critical points inside the narrow slot
dip below $1$. Pure linear interpolation cannot reallocate surplus outer-arc
margin to lift the slot floor.

\subsection{Envelope-Theorem Barycentric Minimax Optimization}
Let $z_0, \dots, z_{M-1}$ be the $M$ Fej\'er nodes on $\partial\Omega$ and
$Q(z) = \prod_{m=0}^{M-1} (z - z_m)$. Every polynomial $P \in \C_{M-1}[z]$ has a
unique barycentric Fej\'er--Lagrange representation
\[
  P(z) = Q(z)\sum_{m=0}^{M-1} \frac{\alpha_m}{z - z_m},
  \qquad f_m := \alpha_m Q'(z_m) = P(z_m).
\]
Crucially, because the Lebesgue constant of Fej\'er nodes on an analytic Jordan
curve grows only as $O(\log M)$, parameterizing $P$ by its $O(1)$-scaled nodal
values $f = (f_0, \dots, f_{M-1}) \in \C^M$ together with the free anchor
$a \in \C$ yields a well-conditioned $2M + 2$ real variable parameterization of
the entire polynomial space $\C_{M-1}[z] \times \C$ (whereas raw weights
$\alpha_m$ vary by five orders of magnitude between the slot and outer arc).

Let $c_1, \dots, c_{M-2}$ be the simple critical points of $P$ ($P'(c_j) = 0$)
and $r_j := |(P(c_j) - P(a)) / ((c_j - a) P'(a))|$. We maximize the smooth
LogSumExp lower-envelope surrogate
\[
  J_\beta(f, a) := -\frac{1}{\beta}\log \sum_{j=1}^{M-2} \exp\!\bigl(-\beta \log r_j\bigr).
\]
Because $P'(c_j) = 0$, the implicit motion $\delta c_j = -\delta P'(c_j)/P''(c_j)$
cancels in $\delta P(c_j) = (\delta P)(c_j) + P'(c_j)\delta c_j = (\delta P)(c_j)$,
leaving an exact closed-form Envelope-Theorem gradient with respect to $(P, a)$:
\begin{equation}\label{eq:envelope}
  \delta \log r_j
  = \RePart\!\left[
      \frac{\delta P(c_j) - \delta P(a)}{P(c_j) - P(a)}
      + \frac{\delta P'(c_j)}{(c_j - a)\, P''(c_j)}
      - \frac{\delta P'(a)}{P'(a)}
    \right],
\end{equation}
plus the direct derivative with respect to $a$. Directly running full nodal
L-BFGS from the initial interpolant creates high-frequency oscillations and
stalls in local traps. Instead, a two-phase schedule cleanly crosses the
macroscopic saddle:
\begin{enumerate}
  \item \textbf{Phase 1 (Smooth Fourier Saddle Crossing, $K = 100$ modes):}
  Restrict updates $\delta f$ to the lowest $K$ discrete Fourier modes on the
  node index (normalized by $M^{-1/2}$ so the anchor $a$ moves at matched scale)
  at $\beta = 400 \to 1{,}500$ without resetting state between stages. This
  smoothly shifts the anchor from $a \approx 0.720 - 0.380\,i$ to
  $a \approx 0.699 - 0.337\,i$, lifting the entire slot floor at the expense of
  surplus outer-arc margin.
  \item \textbf{Phase 2 (Full Nodal L-BFGS Equalization):}
  Unfreeze all $2M + 2$ nodal variables and run L-BFGS with increasing sharpness
  $\beta = 1{,}500 \to 5{,}000 \to 15{,}000 \to 50{,}000$, tracking all $M - 2$
  critical points at every accepted line-search step. On $M = 4{,}000$, this
  immediately produced a degree-$3{,}999$ counterexample
  ($\min r_j \approx 1.0000617$) in $31.9\text{ s}$.
\end{enumerate}

\subsection{Iterative Bootstrap Degree Compression and Maehly Deflation}
To push well below $d = 3{,}999$, we compress each smooth record polynomial
$P_{\mathrm{prev}}$ of degree $M_{\mathrm{prev}} - 1$ onto a smaller Fej\'er grid
$z_0, \dots, z_{M-1}$ by setting $f_m^{(0)} := P_{\mathrm{prev}}(z_m)$, locating
all $M - 2$ critical points of the resampled polynomial, and running Phase~2
L-BFGS equalization. To guarantee $100\%$ root recovery ($M - 2 / M - 2$) at
every compression step even when high-$\beta$ ripples appear, we pair radial
Newton starts with Maehly's deflated Newton iteration,
\[
  \frac{(P'/\Pi)'}{P'/\Pi}(z) = \frac{P''(z)}{P'(z)} - \sum_{i \in \mathcal{R}} \frac{1}{z - r_i},
\]
which systematically repels iterates away from already-discovered roots
$\mathcal{R}$. Stepping progressively through
\[
  d = 3{,}999 \;\longrightarrow\; 3{,}499 \;\longrightarrow\; 3{,}199 \;\longrightarrow\; \mathbf{2{,}799}
\]
yielded our degree-$2{,}799$ polynomial with all $2{,}798$ critical values
$\ge 1.00025213$ at $a \approx 0.698886 - 0.337404\,i$. At $d = 2{,}499$, L-BFGS
equalized both the slot and outer rim at $\min r_j \approx 0.99986644 < 1$,
indicating that $d = 2{,}799$ sits close to the minimax floor of this domain
family while retaining a wide $2.52 \times 10^{-4}$ safety margin for compact
integer rounding.

\subsection{Exact Gaussian-Integer Conversion and Clean Dyadic Anchor}
Using a bit-reversed parallel Karatsuba product tree at $1{,}280$-bit internal
fixed-point precision (\texttt{out\_bits = 256}), the barycentric polynomial is
converted in $3.88\text{ s}$ into exact centered-$z$ monomial form
$P(z) = \sum_{k=0}^{2799} c_k z^k \in \Z[i][z]$ ($1.3\text{ MB}$, maximum coefficient
bit-length $1{,}011$, leading coefficient bit-length $558$). Finally, rounding
the anchor $a$ to the nearest $15$-bit dyadic Gaussian rational
$a = (22901 - 11056\,i)/2^{15}$ shifts $a$ by $< 3.8 \times 10^{-6}$ and
changes $\min r_j$ by less than $5 \times 10^{-8}$.

\section{Two-Stage Zero-Float-Assumption Certification via \texttt{pandroots}}

The standalone Rust certifier \texttt{src/bin/certify\_smale.rs} proves
Theorem~\ref{thm:smale2799} unconditionally without relying on floating-point
arithmetic in any step of the proof:

\begin{enumerate}
  \item \textbf{Exact Degree Verification:}
  Reads $P(z) \in \Z[i][z]$ and the exact integer triple
  $(22901, -11056, 15)$ defining $a \in \Q(i)$ directly from
  \texttt{data/champion\_deg2799.gpoly} and verifies $c_{2799} \neq 0$,
  establishing $\deg(P') = 2{,}798$.

  \item \textbf{Stage 1 --- Exact Critical-Point Certification via
  \texttt{pandroots v0.31.0}~\cite{pandroots}:}
  Constructs the exact Gaussian-integer derivative
  $P'(z) = \sum_{k=1}^{2799} k c_k z^{k-1} \in \Z[i][z]$ and invokes
  \texttt{pandroots::certify\_roots\_exact} with
  \texttt{target\_radius\_bits = 795} (chosen automatically so that
  $M_1(96/100) \cdot 2^{-795} \le |P'(a)|_{\mathrm{lo}} \cdot 2^{-130}$).
  Internally, \texttt{pandroots} polishes the untrusted seeds and then executes
  a zero-float-assumption proof using word-aligned \texttt{u64} limb ball
  arithmetic ($960$ fractional bits):
  \begin{itemize}
    \item \emph{Exact Rouch\'e's theorem:} verifies on every disc
    $D(z_j, r_j)$ ($r_j = 2^{e_j} \le 2^{-795}$) that
    $|P'(z_j)|_{\mathrm{hi}} + \frac{1}{2} M_3 r_j^2 < |P''(z_j)|_{\mathrm{lo}}\, r_j$,
    proving that $D(z_j, r_j)$ contains a unique simple zero of $P'$.
    \item \emph{Exact pairwise disjointness:} an exact integer sweep-line check
    verifies $|z_i - z_j|^2 > (r_i + r_j)^2$ for all $1 \le i < j \le 2{,}798$.
  \end{itemize}
  All $2{,}798 / 2{,}798$ discs pass in $8.56\text{ s}$ ($326.8\text{ roots/s}$).
  By the Fundamental Theorem of Algebra, these $2{,}798$ pairwise-disjoint discs
  contain \emph{every} critical point of $P$ in $\C$, each with multiplicity $1$.

  \item \textbf{Stage 2 --- Exact BigInt Smale Ratio Enclosure:}
  Consuming the exact dyadic centers $z_j$ and power-of-two radii
  $r_j = 2^{e_j}$ returned by \texttt{pandroots} with \emph{zero center
  rounding} at fixed-point precision $F = 1{,}024$ fractional bits:
  \begin{itemize}
    \item Verifies $|z_j|_{\mathrm{hi}} + r_j \le R := 96/100$ in exact integer
    arithmetic for all $2{,}798$ discs.
    \item Computes the exact outward-rounded global derivative bound
    $M_1(R) \ge \sum_{k=1}^{2799} k |c_k| R^{k-1}$, which upper-bounds
    $\sup_{|w| \le R} |P'(w)|$. By convexity of $\overline{D(0, R)}$, every
    $w \in D(z_j, r_j)$ satisfies $|P(w) - P(z_j)| \le M_1(R)\, r_j$.
    \item Evaluates $P(a), P'(a)$, and $P(z_j)$ via exact integer Horner
    arithmetic with explicit ulp rounding bounds (since right-shifting each
    component of a Karatsuba complex product is arithmetic floor division with
    componentwise error in $[0, 1)$ ulp, the Euclidean error per multiplication
    is strictly $< \sqrt{2} < 2$ ulps, propagated at each step as
    $e_p \leftarrow \lfloor e_p |z|_{\mathrm{hi}} / 2^F \rfloor + 3$).
    \item Combines the evaluation and Lipschitz bounds via exact integer
    division to certify $|(P(w) - P(a))/((w - a) P'(a))| > 1$ uniformly for all
    $w \in D(z_j, r_j)$ across all $2{,}798 / 2{,}798$ discs in $1.85\text{ s}$,
    yielding the $25$-digit enclosure~\eqref{eq:certified-min} attained at
    critical point $\#2037$ ($z^* \approx 0.575825497716 - 0.126546779359\,i$).
  \end{itemize}

  \item \textbf{Independent $700$-Digit Python \texttt{mpmath} Cross-Check:}
  The independent script \texttt{verify\_independent.py} audits the exact
  polynomial at $700$ decimal digits on critical point $\#2037$, its $7$ nearest
  neighbors, the $4$ outermost critical points, and $8$ strided boundary roots,
  confirming $|P'(z^*)|/|P'(a)| = 1.03 \times 10^{-690}$ and
  \[
    \left|\frac{P(z^*) - P(a)}{(z^* - a)\, P'(a)}\right|
    = 1.000252081391793447735191684\ldots,
  \]
  agreeing with all $25$ certified digits of~\eqref{eq:certified-min}.
\end{enumerate}

\section{Reproducibility and Artifact Integrity}

The complete source code, exact polynomial, seed file, independent verifier,
and automated GitHub Actions CI workflow (running the full exact
\texttt{pandroots v0.31.0} certification and $700$-digit \texttt{mpmath} audit
on every commit) are publicly archived at
\url{https://github.com/ivan-fr/smale-deg2799}:
\begin{verbatim}
cargo run --release --bin certify_smale
python3 verify_independent.py data/champion_deg2799.gpoly \
    data/champion_deg2799.seeds 2037
\end{verbatim}

\begin{table}[h]
\centering
\footnotesize
\setlength{\tabcolsep}{4pt}
\caption{SHA-256 checksums of the archived degree-$2{,}799$ counterexample artifacts.}
\begin{tabular}{@{}lcl@{}}
\toprule
File & Size & SHA-256 \\
\midrule
\texttt{champion\_deg2799.gpoly} & $1.3\text{ MB}$ & \texttt{\scriptsize 5d1ad7a2509a488934031144d6c2bee5ab36d027bcee392087657a1b117dbfdc} \\
\texttt{champion\_deg2799.seeds} & $128\text{ KB}$ & \texttt{\scriptsize e6fd4dd2a528f58d0b0d690e1401687bbf864cfd6489860841c975ac58045e4a} \\
\bottomrule
\end{tabular}
\end{table}

\begin{thebibliography}{9}

\bibitem{AC26}
T.~Adamczewski and M.~J.~Colbrook,
\emph{From boundary geometry to counterexamples in Smale's mean value problem},
arXiv:2610.36158 (2026).

\bibitem{pandroots}
I.~Besevic,
\emph{\texttt{pandroots}: Certified univariate polynomial root-finder in Rust (v0.31.0)},
\url{https://crates.io/crates/pandroots} (2026).

\bibitem{Besevic7299}
I.~Besevic,
\emph{A machine-certified degree-7,299 counterexample to Smale's mean value conjecture},
Zenodo, \url{https://doi.org/10.5281/zenodo.24512942} (November 24, 2026).

\bibitem{Goel2026}
A.~Goel,
\emph{The sharp constant in Smale's mean value problem},
Preprint (2026).

\bibitem{Smale81}
S.~Smale,
\emph{The fundamental theorem of algebra and complexity theory},
Bull.\ Amer.\ Math.\ Soc.\ (N.S.) \textbf{4} (1981), no.~1, 1--36.

\end{thebibliography}

\end{document}
